\originalchapter{守恒约束与非平衡表面增长}
\label{ch:growth}
\source{8.334 L26：平衡的一般尺度不变性与开放系统增长。}
\originalsection{无临界调参的长程结构}
临界长相关常需调温度到特殊点，但守恒、连续对称或约束也可使长波无质量。不可压缩流体$\nabla\cdot u=0$的均衡动量能$F=\rho\int u^2/2$，允许子空间为横向，故$\avg{u_i(q)u_j(-q)}=(\kb T/\rho)(\delta_{ij}-q_iq_j/q^2)$。尽管振幅无$q$发散，投影在$q=0$方向非解析，实空间含非局域偶极型约束相关。不能仅以“谱有限”判断相关一定纯局域。

这些约束性尺度结构不同于标量液气临界的温度相关方向。开放系统的持续驱动、守恒噪声及模式耦合又可产生新的无临界调参标度，必须从其实际随机方程与对称判断，不能以平衡配分函数自动描述。
\originalsection{Edwards--Wilkinson模型}
表面高度$h(x,t)$相对平均面写$\partial_th=\nu\nabla^2h+\eta$，$\avg{\eta(x,t)\eta(x',t')}=2D\delta^d(x-x')\delta(t-t')$。减掉$q=0$平均高度增长后，每非零模为OU过程，稳态$S_h(q)=D/(\nu q^2)$。

粗糙度$W^2=L^{-d}\int\avg{[h(x)-\bar h]^2}\dd x\sim\int_{1/L}^{1/a}\dd^dq/q^2$。$d<2$给$W\sim L^\alpha$、$\alpha=(2-d)/2$，$d=2$为$W^2\sim\log L$，$d>2$红外有限。动态模时间$\sim q^{-2}$，$z=2$；未饱和时$W\sim t^\beta$，$\beta=\alpha/z=(2-d)/4$。这些指数从精确Gaussian过程来，适于检查模拟和噪声归一。
\originalsection{KPZ非线性与倾斜对称}
若沿表面法向沉积，转成垂直高度速度会出现$(\nabla h)^2$项，KPZ方程为$\partial_th=\nu\nabla^2h+\lambda(\nabla h)^2/2+\eta$。相对平面倾斜$h\to h+v\cdot x$可由坐标漂移$x\to x+\lambda vt$与常数速度补偿，给Galilean倾斜对称。若RG固定点保持该对称且$\lambda$有限，尺度要求$\alpha+z=2$，不是独立地把两个指数都设成Gaussian值。

在一维周期边界，稳态梯度分布仍为$P[h]\propto\exp[-\nu\int(\partial_xh)^2/(2D)]$。Fokker--Planck中非线性流的额外贡献经泛函分部后正比$\int(\partial_xh)^2\partial_x^2h\dd x=\frac13\int\partial_x[(\partial_xh)^3]\dd x=0$。因此$\alpha=1/2$；结合倾斜关系得$z=3/2,\beta=1/3$，不同于EW的$z=2$。这项一维稳态相消一般不在高维成立，也不证明所有噪声或边界下都同属KPZ。
\originalsection{Burgers与Hopf--Cole连接}
令斜率$u=\partial_xh$，得到随机Burgers方程$\partial_tu=\nu\partial_x^2u+\lambda u\partial_xu+\partial_x\eta$，高度不守恒而斜率总积分在周期边界守恒。另一变换$Z=\exp[\lambda h/(2\nu)]$把非线性梯度抵消，形式上$\partial_tZ=\nu\nabla^2Z+\lambda\eta Z/(2\nu)$，是乘性随机热方程。白噪声下需固定Itô/Stratonovich与正规化常数，不能把光滑链式法则直接当作无紫外问题的精确随机恒等式。

随机热方程的路径积分可解释为随机环境中的有向聚合物配分，$h\propto\log Z$是其自由能。这连接非平衡增长与无序多路径问题；下一章讨论无序平均为何不能只把随机势换成均值。
